
The huashen218 corpus is a systematic bibliography of bidirectional human-AI alignment research spanning ML/NLP and HCI communities, assembled by Shen et al. [2024]. Whether these communities cite each other's alignment work symmetrically is an open empirical question with structural implications for the field. This paper builds and validates the first Semantic Scholar Academic Graph (S2AG)-based bibliometric pipeline for within-corpus directed citation analysis of this corpus. In doing so, a significant methodological finding emerges: standard venue-string classification fails on interdisciplinary corpora, labeling only 12\% of resolved papers from the huashen218 reading list, compared to 100\% coverage achieved with FoS-primary classification --- a method that uses S2AG \texttt{fieldsOfStudy} tags rather than venue name strings. In the 33 papers resolved from the public reading list, HCI alignment papers directed all within-corpus outgoing citations to ML/NLP work, while ML/NLP alignment papers allocated only 10.6\% of outgoing citations to HCI, yielding a directed citation ratio of approximately 0.107. Both pre-registered gate thresholds --- resolution coverage $\geq$ 70\% and $\geq$ 30 cross-group edges --- were not met (67.3\% and 9 edges, respectively), so no statistical tests were executed and all three quantitative predictions are classified as INCONCLUSIVE rather than REFUTED. The validated, fully cached pipeline is available for re-execution once the full corpus is reconstructed programmatically.

